\documentclass[a4paper,11pt]{article}
\usepackage[utf8]{inputenc}
\usepackage[T1]{fontenc}
\usepackage[french]{babel}
\usepackage{geometry}
\geometry{left=2cm,right=2cm,top=1.5cm,bottom=1.5cm}
\usepackage{hyperref}
\usepackage{enumitem}

\begin{document}

\begin{flushleft}
    \textbf{AISSAOUI Ismail} \\
    Ingénieur en Intelligence Artificielle \& Data Science \\
    Chlef, Algérie \\
    Tél : +213 660 70 77 96 \\
    Email : \href{mailto:ismail.aissaoui.pro@gmail.com}{ismail.aissaoui.pro@gmail.com} \\
    Portfolio : \href{https://ismail-aissaoui.vercel.app}{ismail-aissaoui.vercel.app}
\end{flushleft}

\begin{flushright}
    À l'attention de \\
    \textbf{M. Baptiste Caramiaux} \\
    \textbf{M. Gonzalo Ramos} \\
    Équipe HCI Sorbonne, ISIR \\
    Sorbonne Université \\
    Fait le \today
\end{flushright}

\vspace{0.4cm}

\noindent \textbf{Objet : Candidature au sujet de thèse : « Nouveaux paradigmes d’interaction pour agir avec et à travers l’IA générative »}

\vspace{0.4cm}

Monsieur Caramiaux, Monsieur Ramos,

\medskip

Ingénieur en Intelligence Artificielle diplômé de l'ESI-SBA, c'est avec une immense motivation que je vous soumets ma candidature pour ce projet doctoral. Votre proposition, qui traite l'IA générative non pas comme un simple outil de réponse, mais comme un \textbf{matériau et un médium} pour l'action médiée, définit exactement le territoire de recherche que je souhaite explorer : le dépassement du paradigme restrictif du « prompting » vers des architectures interactives fluides.

Mon parcours m'a déjà permis de confronter les défis fondamentaux cités dans votre offre, notamment le franchissement des \textit{fossés d'exécution et d'évaluation} (Norman, 1985). À travers mon projet \textit{Geo-RAG}, j'ai architecturé un système où la GenAI sert de substrat pour l'exploration de connaissances spatiales, en mettant l'accent sur le maintien du contexte et l'expression d'intentions complexes. De même, mon travail sur \textit{Recognizini} a été une première incursion dans la redistribution de l'agentivité via des boucles de feedback itératives, permettant à l'utilisateur d'agir « à travers » le système pour affiner l'inférence en temps réel.

Le projet de thèse s'articule autour de quatre questions de recherche cruciales auxquelles je suis impatient de répondre. Ma maîtrise technique des modèles de langage (Transformers, architectures agentiques) et mon expertise en prototypage frontend (React, Framer Motion) me permettront de concevoir et d'implémenter les architectures systèmes exposant les états de raisonnement intermédiaires que vous envisagez. Je suis convaincu que mon profil hybride est idéal pour passer de la formalisation théorique à la création de prototypes exemplaires ciblant les conférences majeures comme CHI, UIST ou IUI.

Intégrer l'équipe ACIDE à l'ISIR serait pour moi l'opportunité de participer à ce recadrage profond de l'IHM, en développant de nouveaux langages d'interaction qui préservent le contrôle et la responsabilité humaine dans un monde de plus en plus médiatisé par l'IA.

Je vous remercie de l'attention que vous porterez à ma candidature et reste à votre entière disposition pour un entretien.

Je vous prie d'agréer, Monsieur Caramiaux, Monsieur Ramos, l'expression de mes salutations respectueuses.

\vspace{0.4cm}

\begin{flushleft}
    \textbf{AISSAOUI Ismail}
\end{flushleft}

\end{document}
